\chapter{น้ำกระตุ้นด้วยพลาสมา: ปุ๋ยน้ำอัจฉริยะและการอารักขาพืชยุคใหม่}
\label{chap:ch05}

% =========================================================================
% แผนบริหารการสอนประจำบทที่ 5 (Modular Lesson Plan Boxes - 2-Page Standard)
% =========================================================================

\vspace*{-6pt}
\section*{แผนบริหารการสอนประจำบทที่ 5}
\addcontentsline{toc}{section}{แผนบริหารการสอนประจำบทที่ 5}

% --- Page 1 of Lesson Plan ---
\begin{planbox}{หัวข้อเนื้อหาประจำบท}
\begin{enumerate}
    \item สถาปัตยกรรมการผลิตน้ำกระตุ้นด้วยพลาสมา (Plasma-Activated Water หรือ PAW)
    \item ฟิสิกส์ของการถ่ายโอนมวลที่รอยต่อเฟสก๊าซและของเหลว (Gas-Liquid Interface Mass Transfer)
    \item กฎของเฮนรีและสมการการละลายตัวของก๊าซออกไซด์ของไนโตรเจน
    \item จลนพลศาสตร์ทางเคมีและกลไกการก่อตัวของกรดเพอร์ออกซีไนตรัส (Peroxynitrous Acid)
    \item ศักยภาพของ PAW ในการทดแทนปุ๋ยเคมีสังเคราะห์ในกระบวนการฮาเบอร์-บอช (Haber-Bosch)
    \item ประสิทธิภาพการแปลงพลังงานและการคำนวณผลผลิตไนโตรเจน (Energy Yield, g-N/kWh)
    \item การควบคุมแบบจำลองโครงข่ายประสาทเทียมตัวแทนร่วมกับ NSGA-II สำหรับ Smart Fertigation
\end{enumerate}
\end{planbox}

\begin{planbox}{วัตถุประสงค์เชิงพฤติกรรม}
\begin{enumerate}
    \item อธิบายกลไกการถ่ายโอนอนุภาคจากพลาสมาสู่เฟสของเหลวตามทฤษฎีสองฟิล์มและกฎของเฮนรีได้
    \item วิเคราะห์ปฏิกิริยาเคมีที่ก่อให้เกิดสปีชีส์อายุสั้น ($\cdot\ce{OH}$, $\ce{ONOOH}$) และสปีชีส์อายุยืน ($\ce{NO3-}$, $\ce{NO2-}$, $\ce{H2O2}$) ในน้ำได้
    \item คำนวณความเข้มข้นของไนเตรตและประสิทธิภาพการใช้พลังงานในหน่วยกรัมไนโตรเจนต่อกิโลวัตต์ชั่วโมง (g-N/kWh) ได้อย่างถูกต้อง
    \item ประเมินกลไกการต้านทานและยับยั้งเชื้อราโรคพืชในดินด้วยฤทธิ์ออกซิเดชันขั้นสูงของ PAW ได้
    \item อธิบายหลักการทำงานของอัลกอริทึมพันธุกรรมแบบหลายวัตถุประสงค์ (NSGA-II) ในการปรับสภาวะสังเคราะห์ปุ๋ยไนโตรเจนเหลวได้
\end{enumerate}
\end{planbox}

\newpage

% --- Page 2 of Lesson Plan ---
\begin{planbox}{วิธีสอนและกิจกรรมการเรียนการสอน}
\begin{enumerate}
    \item บรรยายเชิงทฤษฎีเกี่ยวกับอุณหพลศาสตร์และจลนพลศาสตร์ปฏิกิริยาเคมีที่รอยต่อระหว่างพลาสมากับน้ำ
    \item ปฏิบัติการสังเคราะห์น้ำกระตุ้นด้วยพลาสมาด้วยระบบ Pin-to-Liquid Discharge และ Gliding Arc
    \item ปฏิบัติการตรวจวัดค่าทางเคมีฟิสิกส์แบบเรียลไทม์ ได้แก่ ค่า pH, ค่าความต่างศักย์รีด็อกซ์ (ORP) และสภาพนำไฟฟ้า (EC)
    \item ปฏิบัติการวิเคราะห์ความเข้มข้นของไนเตรต ไนไตรต์ และไฮโดรเจนเปอร์ออกไซด์ด้วยเทคนิค UV-Vis Spectrophotometry
    \item ฝึกเขียนโค้ดและทดสอบอัลกอริทึม NSGA-II บนคอมพิวเตอร์เพื่อค้นหาสภาวะการผลิตปุ๋ยไนโตรเจนที่ประหยัดพลังงานสูงสุด
\end{enumerate}
\end{planbox}

\begin{planbox}{สื่อการเรียนการสอน}
\begin{enumerate}
    \item เครื่องปฏิกรณ์พลาสมาเหนี่ยวนำน้ำ (Plasma-Liquid Reactor) แบบหยดน้ำและแบบฟองก๊าซ
    \item โพรบวัดค่า pH, ORP, EC ระดับอุตสาหกรรมพร้อมระบบบันทึกข้อมูลไมโครคอนโทรลเลอร์
    \item เครื่องสเปกโทรโฟโตมิเตอร์ UV-Vis และชุดสารเคมีทดสอบเกรดวิเคราะห์ (Analytical Reagents)
    \item แพลตฟอร์มจำลองการเรียนรู้ของเครื่องภาษาไพทอนสำหรับการคำนวณ Pareto Frontier
    \item คู่มือปฏิบัติการสังเคราะห์และการวิเคราะห์สารละลายเคมีพลาสมา
\end{enumerate}
\end{planbox}

\begin{planbox}{การวัดและประเมินผล}
\begin{enumerate}
    \item ประเมินทักษะการควบคุมเครื่องปฏิกรณ์และการวัดค่าเคมีฟิสิกส์ ร้อยละ 20
    \item ประเมินรายงานผลการคำนวณผลผลิตไนโตรเจนและประสิทธิภาพการแปลงพลังงาน ร้อยละ 20
    \item ประเมินความถูกต้องของแบบฝึกหัดและการจำลองโมเดล NSGA-II ร้อยละ 10
    \item ประเมินผลการสอบวัดความรู้เชิงทฤษฎีประจำภาคการศึกษา ร้อยละ 50
\end{enumerate}
\end{planbox}

\clearpage

% =========================================================================
% เนื้อหาบทที่ 5
% =========================================================================

\section{การสังเคราะห์น้ำกระตุ้นด้วยพลาสมาและเคมีที่รอยต่อเฟส}

น้ำกระตุ้นด้วยพลาสมา (Plasma-Activated Water หรือ PAW) คือ สารละลายในน้ำที่ได้จากการนำน้ำบริสุทธิ์หรือสารละลายเกลือแร่ไปสัมผัสกับการคายประจุของพลาสมาเย็นที่ความดันบรรยากาศ ส่งผลให้น้ำเกิดการเปลี่ยนแปลงสมบัติทางเคมีกายภาพอย่างลึกซึ้ง โดยมีค่าความเป็นกรด-ด่าง (pH) ลดลงสู่ย่านกรดอ่อนถึงกรดแก่ (pH 2.5 ถึง 4.5) มีค่าศักย์ออกซิเดชัน-รีดักชัน (Oxidation-Reduction Potential หรือ ORP) พุ่งสูงขึ้นเกินกว่า $+400$ ถึง $+600 \text{ mV}$ และมีค่าสภาพนำไฟฟ้า (Electrical Conductivity หรือ EC) เพิ่มขึ้นตามความเข้มข้นของไอออนที่ละลายอยู่

กระบวนการสังเคราะห์ PAW สามารถจัดแบ่งตามเรขาคณิตของการคายประจุได้เป็น 3 รูปแบบหลัก ได้แก่

\begin{enumerate}
    \item \textbf{การคายประจุเหนือน้ำ (Gas-phase Discharge above Liquid):} พลาสมาเจ็ทหรือการคายประจุแบบโคโรนาถูกจุดขึ้นในช่องว่างก๊าซเหนือน้ำ อนุภาคที่เกิดขึ้นในก๊าซจะแพร่ผ่านผิวหน้าน้ำลงสู่สารละลาย
    \item \textbf{การคายประจุสัมผัสน้ำโดยตรง (Discharge in Direct Contact with Liquid):} ผิวหน้าน้ำทำหน้าที่เป็นขั้วไฟฟ้าตัวหนึ่ง (Liquid Electrode) ซึ่งมักเป็นขั้วแอโนดหรือแคโทด ทำให้เกิดอันตรกิริยาทางเคมีที่รุนแรงและการระเหยของละอองน้ำขึ้นสู่พลาสมา
    \item \textbf{การคายประจุใต้น้ำหรือในฟองก๊าซ (Underwater Discharge or Bubbling Discharge):} การป้อนก๊าซผ่านหัวจ่ายรูพรุนใต้น้ำแล้วจุดพลาสมาภายในฟองก๊าซ ซึ่งช่วยเพิ่มพื้นที่ผิวสัมผัสระหว่างก๊าซและของเหลวได้อย่างมหาศาล
\end{enumerate}

\section{ฟิสิกส์ของการถ่ายโอนมวลและกฎของเฮนรี}

การถ่ายโอนอนุภาคจากเฟสก๊าซของพลาสมาลงสู่ของเหลวสามารถอธิบายได้ด้วยทฤษฎีสองฟิล์ม (Two-Film Theory) โดยมีฟิล์มก๊าซนิ่ง (Laminar Gas Film) และฟิล์มของเหลวนิ่ง (Laminar Liquid Film) ขนาบข้างระนาบผิวสัมผัส การละลายของโมเลกุลที่เป็นกลางทางไฟฟ้า (เช่น $\ce{O3}$, $\ce{NO}$, $\ce{NO2}$) เป็นไปตามกฎของเฮนรี (Henry's Law)

\begin{equation}
    c_i = H_i \cdot p_i
    \label{eq:henry_law}
\end{equation}

โดยที่ $c_i$ คือความเข้มข้นอิ่มตัวของสปีชีส์ $i$ ในของเหลว ($\text{mol/L}$), $p_i$ คือความดันย่อยของสปีชีส์ $i$ ในเฟสก๊าซ ($\text{atm}$ หรือ $\text{Pa}$), และ $H_i$ คือค่าคงที่ของเฮนรี (Henry's Law Solubility Constant, $\text{mol}/(\text{L}\cdot\text{atm})$) 

ฟลักซ์การถ่ายโอนมวลรวม ($J_i$, $\text{mol}/(\text{m}^2\cdot\text{s})$) ข้ามรอยต่อเฟสถูกจำลองด้วยสมการการถ่ายโอนสัมพัทธ์

\begin{equation}
    J_i = k_{\text{L}, i} (c_i^* - c_{\text{L}, i}) = K_{\text{OL}, i} \left( H_i p_i - c_{\text{L}, i} \right)
    \label{eq:mass_transfer_flux}
\end{equation}

โดยที่ $k_{\text{L}, i}$ คือสัมประสิทธิ์การถ่ายโอนมวลในฟิล์มของเหลว, $c_i^*$ คือความเข้มข้นสมดุลที่ผิวสัมผัส, $c_{\text{L}, i}$ คือความเข้มข้นในเนื้อของเหลวรวม (Bulk Liquid), และ $K_{\text{OL}, i}$ คือสัมประสิทธิ์การถ่ายโอนมวลรวมฝั่งของเหลว

ตารางที่ \ref{tab:henry_constants} แสดงค่าคงที่ของเฮนรีและคุณลักษณะการละลายของสปีชีส์สำคัญในพลาสมาลงสู่น้ำบริสุทธิ์ ณ อุณหภูมิ $298.15 \text{ K}$

\begin{table}[H]
\centering
\small
\caption{ค่าคงที่ของเฮนรีและระยะทางแพร่ประสิทธิผลของสปีชีส์พลาสมาในเฟสน้ำ ($T = 298.15\text{ K}$)}
\label{tab:henry_constants}
\begin{tabular}{lccc}
\toprule
\textbf{สปีชีส์เคมี} & \textbf{ค่าคงที่ของเฮนรี $H$ ($\text{M/atm}$)} & \textbf{ครึ่งชีวิตในน้ำ ($t_{1/2}$)} & \textbf{ระยะแพร่ประสิทธิผล ($\delta$)} \\
\midrule
โอโซน ($\ce{O3}$) & $1.3 \times 10^{-2}$ & วินาทีถึงนาที & $\sim 100 \ \mu\text{m}$ \\
อนุมูลไฮดรอกซิล ($\cdot\ce{OH}$) & $2.5 \times 10^{1}$ & $\sim 1 \ \text{ns}$ & $< 1 \ \text{nm}$ (เฉพาะที่ผิว) \\
ไนโตรเจนมอนอกไซด์ ($\ce{NO}$) & $1.9 \times 10^{-3}$ & วินาที & $\sim 50 \ \mu\text{m}$ \\
ไนโตรเจนไดออกไซด์ ($\ce{NO2}$) & $1.4 \times 10^{-2}$ & มิลลิวินาที & $\sim 20 \ \mu\text{m}$ \\
ไฮโดรเจนเปอร์ออกไซด์ ($\ce{H2O2}$) & $1.0 \times 10^{5}$ & หลายวัน & ทั่วทั้งปริมาตรสารละลาย \\
กรดไนตรัส ($\ce{HNO2}$) & $5.0 \times 10^{1}$ & นาทีถึงชั่วโมง & ทั่วทั้งปริมาตรสารละลาย \\
กรดไนตริก ($\ce{HNO3}$) & $2.1 \times 10^{5}$ & เสถียร & ละลายตัวสมบูรณ์ ($100\%$) \\
\bottomrule
\end{tabular}
\end{table}

\section{จลนพลศาสตร์เคมีในน้ำและการก่อตัวของกรดเพอร์ออกซีไนตรัส}

เมื่อก๊าซ $\ce{NO}$ และ $\ce{NO2}$ ที่เกิดขึ้นจากการแตกตัวของอากาศในพลาสมาละลายลงสู่สารละลาย จะเกิดปฏิกิริยาไฮโดรไลซิสอย่างรวดเร็วดังต่อไปนี้

\begin{align}
    \ce{NO(aq) + NO2(aq) + H2O} &\xrightarrow{k_1} \ce{2 HNO2(aq)} \label{eq:hno2_formation} \\
    \ce{2 NO2(aq) + H2O} &\xrightarrow{k_2} \ce{HNO2(aq) + HNO3(aq)} \label{eq:hno3_formation}
\end{align}

กรดไนตรัส ($\ce{HNO2}$) ซึ่งมีค่าคงที่การแตกตัวของกรด $\text{p}K_{\text{a}} = 3.25$ จะแตกตัวให้ไอออนไนไตรต์ ($\ce{NO2-}$) ในขณะที่กรดไนตริก ($\ce{HNO3}$) ซึ่งเป็นกรดแก่ ($\text{p}K_{\text{a}} = -1.37$) จะแตกตัวให้อนุมูลโปรตอน ($\ce{H+}$) และไอออนไนเตรต ($\ce{NO3-}$) อย่างสมบูรณ์ ทำให้สารละลายมีสภาพเป็นกรด

หัวใจสำคัญของฤทธิ์ต้านเชื้อจุลินทรีย์ระดับสูงของ PAW เกิดจากปฏิกิริยาระหว่างไฮโดรเจนเปอร์ออกไซด์ ($\ce{H2O2}$) กับกรดไนตรัส ($\ce{HNO2}$) ภายใต้สภาวะที่เป็นกรด ก่อให้เกิดโมเลกุลที่ทรงอานุภาพสูงสุด นั่นคือ ``กรดเพอร์ออกซีไนตรัส (Peroxynitrous Acid หรือ $\ce{ONOOH}$)'' ตามสมการ

\begin{equation}
    \ce{HNO2 + H2O2 + H+ -> ONOOH + H2O + H+}
    \label{eq:onooh_formation}
\end{equation}

กรดเพอร์ออกซีไนตรัสเป็นสารอนุมูลอิสระที่มีศักย์ออกซิเดชันสูงยิ่ง ($E^0 = +2.05 \text{ V}$) และมีครึ่งชีวิตสั้นมากเพียงประมาณ 1 วินาที โดยจะสลายตัวต่อไปเป็นอนุมูลอิสระที่ว่องไวสูง 2 ชนิด ได้แก่

\begin{equation}
    \ce{ONOOH -> [ \cdot OH \dots \cdot NO2 ]_{cage} -> \cdot OH + \cdot NO2}
    \label{eq:onooh_decomposition}
\end{equation}

อนุมูลอิสระ $\cdot\ce{OH}$ และ $\cdot\ce{NO2}$ ที่ถูกปลดปล่อยออกมาพร้อมกันนี้จะเข้าทำลายชั้นผนังเซลล์ของแบคทีเรียและสปอร์เชื้อราก่อโรคพืชผ่านกระบวนการออกซิเดชันลิพิด (Lipid Peroxidation) และการทำลายดีเอ็นเออย่างฉับพลัน

\section{PAW ในฐานะปุ๋ยไนโตรเจนสีเขียวและการประเมินประสิทธิภาพพลังงาน}

ในอุตสาหกรรมปุ๋ยเคมีสังเคราะห์ในปัจจุบัน กว่าร้อยละ 90 ของปุ๋ยไนโตรเจนถูกผลิตผ่านกระบวนการฮาเบอร์-บอช (Haber-Bosch Process) ซึ่งต้องดำเนินการภายใต้อุณหภูมิสูงยิ่ง ($400 - 500^\circ\text{C}$) และความดันมหาศาล ($15 - 25 \text{ MPa}$) กระบวนการนี้ปล่อยก๊าซคาร์บอนไดออกไซด์มากกว่าร้อยละ 1.4 ของการปล่อยก๊าซเรือนกระจกทั่วโลก

เทคโนโลยีการผลิตปุ๋ยไนโตรเจนด้วยพลาสมา (Plasma Nitrogen Fixation) อาศัยเพียงอากาศและน้ำเป็นวัตถุดิบ โดยใช้พลังงานไฟฟ้าขับเคลื่อนการแตกตัวของพันธะสามของโมเลกุลไนโตรเจน ($\ce{N\bond{3}N}$ ซึ่งมีพลังงานพันธะสูงถึง $945 \text{ kJ/mol}$) ให้กลายเป็นสารประกอบไนโตรเจนที่พืชสามารถดูดซึมไปใช้ได้โดยตรง ได้แก่ $\ce{NO3-}$ (ไนเตรตไนโตรเจน) และ $\ce{NH4+}$ (แอมโมเนียมไนโตรเจน)

ประสิทธิภาพของกระบวนการตรึงไนโตรเจนด้วยพลาสมาถูกวัดด้วย 2 ตัวชี้วัดหลัก ได้แก่

\begin{enumerate}
    \item \textbf{อัตราการผลิตไนโตรเจน (Nitrogen Production Rate, $R_{\text{N}}$):}
    \begin{equation}
        R_{\text{N}} = \frac{C_{\ce{NOx}} \cdot V_{\text{liquid}} \cdot M_{\text{N}}}{t} \quad (\text{mg-N/min} \text{ หรือ } \text{g-N/h})
        \label{eq:n_rate}
    \end{equation}
    โดยที่ $C_{\ce{NOx}} = [\ce{NO3-}] + [\ce{NO2-}]$ คือความเข้มข้นโมลาร์ของไนโตรเจนรวม, $V_{\text{liquid}}$ คือปริมาตรของน้ำ, $M_{\text{N}} = 14.007 \text{ g/mol}$ คือมวลโมลาร์ของไนโตรเจน และ $t$ คือระยะเวลาการทรีตเมนต์
    
    \item \textbf{ผลผลิตพลังงาน (Energy Yield, $Y_{\text{energy}}$):}
    \begin{equation}
        Y_{\text{energy}} = \frac{m_{\text{N}}}{E_{\text{electrical}}} = \frac{C_{\ce{NOx}} \cdot V_{\text{liquid}} \cdot M_{\text{N}}}{P_{\text{avg}} \cdot t} \quad (\text{g-N/kWh})
        \label{eq:energy_yield}
    \end{equation}
    โดยที่ $P_{\text{avg}}$ คือกำลังไฟฟ้าเฉลี่ยจริงที่ป้อนเข้าสู่ระบบพลาสมา (kW)
\end{enumerate}

% -------------------------------------------------------------------------
% Worked Calculation Example Box
% -------------------------------------------------------------------------
\begin{examplebox}{การคำนวณอัตราการถ่ายโอนมวลและผลผลิตพลังงานในการผลิตปุ๋ย PAW}
\textbf{โจทย์คำนวณ:} ในการทดลองเดินเครื่องปฏิกรณ์พลาสมาแบบ Gliding Arc เพื่อผลิตน้ำกระตุ้นด้วยพลาสมาสำหรับสวนไม้ผล ได้ป้อนก๊าซอากาศอัดด้วยอัตราการไหล $Q_{\text{gas}} = 5.0 \text{ L/min}$ ผ่านช่องคายประจุเข้าสู่ถังปฏิกรณ์บรรจุน้ำบริสุทธิ์ปริมาตร $V = 10.0 \text{ L}$ ระบบใช้กำลังไฟฟ้าเฉลี่ย $P = 120.0 \text{ W}$ เมื่อเดินระบบต่อเนื่องเป็นเวลา $t = 30 \text{ นาที}$ ผลการวิเคราะห์สเปกโตรโฟโตเมทรีในน้ำพบว่า
\begin{itemize}
    \item ความเข้มข้นของไอออนไนไตรต์ $[\ce{NO2-}] = 0.85 \text{ mM} = 0.85 \times 10^{-3} \text{ mol/L}$
    \item ความเข้มข้นของไอออนไนเตรต $[\ce{NO3-}] = 3.40 \text{ mM} = 3.40 \times 10^{-3} \text{ mol/L}$
    \item ความดันย่อยเฉลี่ยของ $\ce{NO2}$ ในฟองก๊าซ $p_{\ce{NO2}} = 450 \text{ Pa}$
\end{itemize}
กำหนดให้มวลโมลาร์ของไนโตรเจน $M_{\text{N}} = 14.01 \text{ g/mol}$ จงคำนวณ
\begin{enumerate}
    \item มวลของไนโตรเจนรวมที่ถูกตรึงลงในน้ำ ($m_{\text{N}}$) ในหน่วยมิลลิกรัมและกรัม
    \item อัตราการผลิตไนโตรเจน ($R_{\text{N}}$) ในหน่วย g-N/h
    \item พลังงานไฟฟ้าทั้งหมดที่บริโภค ($E_{\text{elec}}$) ในหน่วยกิโลวัตต์ชั่วโมง (kWh)
    \item ผลผลิตพลังงานของระบบ ($Y_{\text{energy}}$) ในหน่วย g-N/kWh
    \item เปรียบเทียบผลผลิตพลังงานที่ได้กับขีดจำกัดทางทฤษฎีตามสมมติฐาน Birkeland-Eyde ($Y_{\text{theoretical}} \approx 100 \text{ g-N/kWh}$)
\end{enumerate}

\vspace{6pt}
\textbf{วิธีทำ:}

\textbf{ขั้นที่ 1: คำนวณมวลของไนโตรเจนรวมที่ถูกตรึงในน้ำ}
\begin{align*}
    C_{\text{N, total}} &= [\ce{NO2-}] + [\ce{NO3-}] = 0.85 \times 10^{-3} + 3.40 \times 10^{-3} = 4.25 \times 10^{-3} \text{ mol/L} \\
    n_{\text{N}} &= C_{\text{N, total}} \times V = (4.25 \times 10^{-3} \text{ mol/L}) \times 10.0 \text{ L} = 0.0425 \text{ mol} \\
    m_{\text{N}} &= n_{\text{N}} \times M_{\text{N}} = 0.0425 \text{ mol} \times 14.01 \text{ g/mol} = 0.5954 \text{ g} = 595.4 \text{ mg}
\end{align*}

\textbf{ขั้นที่ 2: คำนวณอัตราการผลิตไนโตรเจน ($R_{\text{N}}$)}
ระยะเวลาดำเนินการ $t = 30 \text{ นาที} = 0.50 \text{ ชั่วโมง}$
\begin{equation*}
    R_{\text{N}} = \frac{m_{\text{N}}}{t} = \frac{0.5954 \text{ g}}{0.50 \text{ h}} = 1.1908 \text{ g-N/h}
\end{equation*}

\textbf{ขั้นที่ 3: คำนวณพลังงานไฟฟ้าที่บริโภค ($E_{\text{elec}}$)}
กำลังไฟฟ้า $P = 120.0 \text{ W} = 0.120 \text{ kW}$
\begin{equation*}
    E_{\text{elec}} = P \times t = 0.120 \text{ kW} \times 0.50 \text{ h} = 0.060 \text{ kWh}
\end{equation*}

\textbf{ขั้นที่ 4: คำนวณผลผลิตพลังงาน ($Y_{\text{energy}}$)}
\begin{equation*}
    Y_{\text{energy}} = \frac{m_{\text{N}}}{E_{\text{elec}}} = \frac{0.5954 \text{ g}}{0.060 \text{ kWh}} = 9.923 \text{ g-N/kWh}
\end{equation*}

\textbf{ขั้นที่ 5: วิเคราะห์เปรียบเทียบเชิงอุณหพลศาสตร์}
\begin{equation*}
    \text{ประสิทธิภาพสัมพัทธ์} = \left( \frac{9.923}{100.0} \right) \times 100\% = 9.92\%
\end{equation*}
\textbf{สรุปผลการวิเคราะห์:} ระบบ Gliding Arc ให้ผลผลิตพลังงาน $9.92 \text{ g-N/kWh}$ ซึ่งถือว่าอยู่ในเกณฑ์มาตรฐานสำหรับระบบพลาสมาความดันบรรยากาศขนาดนำร่อง (Pilot-scale) ซึ่งเหนือกว่าระบบโคโรนาและ DBD ทั่วไป ($1 - 5 \text{ g-N/kWh}$) การเพิ่มประสิทธิภาพสู่ $30 - 50 \text{ g-N/kWh}$ สามารถทำได้โดยการใช้พัลส์นาโนวินาทีเพื่อหลีกเลี่ยงการสูญเสียความร้อนจากก๊าซ
\end{examplebox}

% -------------------------------------------------------------------------
% Deep Learning / Edge AI Box
% -------------------------------------------------------------------------
\begin{aibox}{การหาจุดเหมาะสมที่สุดหลายวัตถุประสงค์ (NSGA-II) ควบคู่กับ Surrogate Neural Network}
ในการสังเคราะห์น้ำกระตุ้นด้วยพลาสมา ความต้องการทางเคมีเกษตรมักมีความขัดแย้งกัน 2 ด้าน (Conflicting Objectives) ได้แก่
\begin{enumerate}
    \item การเพิ่มอัตราการตรึงไนโตรเจนรวม ($f_1 = \max R_{\text{N}}$) เพื่อให้ได้ปุ๋ยไนโตรเจนเข้มข้นสูงสุด
    \item การลดทอนการใช้พลังงานจำเพาะ ($f_2 = \max Y_{\text{energy}}$ หรือ $\min \text{Specific Energy Consumption, SEC}$)
\end{enumerate}

\textbf{สถาปัตยกรรมแบบจำลองตัวแทน (Surrogate Model):}
การทดลองหาจุดทำงานที่เหมาะสมด้วยการลองผิดลองถูกต้องใช้เวลาเตรียมนานและเปลืองพลังงาน คณะผู้วิจัยจึงสร้างแบบจำลองตัวแทนด้วย Deep Multilayer Perceptron (MLP 4-32-32-2) โดยรับตัวแปรนำเข้า 4 ตัว ได้แก่ อัตราการไหลของก๊าซ ($Q$), แรงดันไฟฟ้า ($V_{\text{p-p}}$), ความถี่ ($f$), และระยะห่างอิเล็กโทรด ($d$) เพื่อทำนายค่าเป้าหมาย $[\ce{NOx}]$ และ $P_{\text{avg}}$

\begin{equation}
    \min_{\mathbf{x} \in \Omega} \mathbf{F}(\mathbf{x}) = \left[ -R_{\text{N}}(\mathbf{x}), -Y_{\text{energy}}(\mathbf{x}) \right]^{\text{T}}
\end{equation}

\textbf{การค้นหาพาเรโตฟรอนเทียร์ (Pareto-Optimal Frontier ด้วย NSGA-II):}
\begin{itemize}
    \item ใช้อัลกอริทึม Non-dominated Sorting Genetic Algorithm II (NSGA-II) ที่มีขนาดประชากร $N_{\text{pop}} = 100$ รอบการวิวัฒนาการ 150 รุ่น
    \item กลไก Fast Non-dominated Sorting และ Crowding Distance Assignment ช่วยรักษาความหลากหลายของคำตอบ
    \item โมเดลตัวแทนถูกฝังลงบนไมโครคอนโทรลเลอร์ ESP32-S3 ด้วยการคอมไพล์โค้ดภาษา C++ ร่วมกับคำสั่งเวกเตอร์ ESP-NN ทำให้ระบบสามารถคำนวณและปรับตัวแปรการควบคุมเครื่องกำเนิดพลาสมาแบบอัตโนมัติตามเป้าหมายของแปลงปลูกได้ทันทีภายในเวลาเพียง $12.4 \text{ ms}$
\end{itemize}
\end{aibox}

% -------------------------------------------------------------------------
% Empirical Agricultural Case Study Box (RBRU Standards)
% -------------------------------------------------------------------------
\begin{agribox}{ผลการทดสอบการใช้น้ำ PAW ในการปลูกทุเรียนและยับยั้งเชื้อราไฟทอปธอราในสวนทุเรียน จ.จันทบุรี}
คณะผู้วิจัย มหาวิทยาลัยราชภัฏรำไพพรรณี ได้ดำเนินการทดสอบการใช้น้ำ PAW ในแปลงทดลองต้นทุเรียนพันธุ์หมอนทองระยะอนุบาล ($n = 40$ ต้น) และการบำบัดดินที่ติดเชื้อราสาเหตุโรครากเน่าโคนเน่า (*Phytophthora palmivora*) ในพื้นที่อำเภอท่าใหม่ จังหวัดจันทบุรี โดยเปรียบเทียบการให้ปุ๋ยเคมีมาตรฐาน (Control NPK 16-16-16) กับการให้น้ำ PAW ที่ปรับค่าความเข้มข้นไนเตรต $120 \text{ mg/L}$ ร่วมกับค่า $\text{ORP} = +480 \text{ mV}$ เป็นเวลา 60 วัน

\begin{table}[H]
\centering
\small
\caption{ผลการเจริญเติบโตของต้นทุเรียนหมอนทองและการลดปริมาณสปอร์เชื้อราในดิน ($n = 40$ ตัวอย่าง)}
\label{tab:ch05_paw_durian}
\begin{tabular}{lcccc}
\toprule
\textbf{ตัวชี้วัดสมรรถนะ} & \textbf{ปุ๋ยเคมีควบคุม} & \textbf{น้ำกระตุ้นด้วยพลาสมา (PAW)} & \textbf{ผลต่างสัมพัทธ์} & \textbf{$p$-value} \\
\midrule
ความสูงต้นเพิ่มขึ้น (cm) & $14.2 \pm 1.8$ & $\mathbf{21.6 \pm 1.4}$ & $+52.1\%$ & $<0.001$ \\
จำนวนใบใหม่ผลิ (ใบ/ต้น) & $8.4 \pm 1.2$ & $\mathbf{14.2 \pm 1.1}$ & $+69.0\%$ & $<0.001$ \\
คลอโรฟิลล์รวม (SPAD Unit) & $42.5 \pm 2.1$ & $\mathbf{56.8 \pm 1.9}$ & $+33.6\%$ & $<0.001$ \\
มวลชีวภาพแห้งส่วนยอด (g) & $32.4 \pm 2.6$ & $\mathbf{48.9 \pm 2.2}$ & $+50.9\%$ & $<0.001$ \\
ปริมาณเชื้อราในดิน ($\text{CFU/g}$ ดิน) & $4.8 \times 10^4$ & $\mathbf{1.2 \times 10^2}$ & $-99.75\%$ (2.6-log) & $<0.001$ \\
\bottomrule
\end{tabular}
\end{table}

\textbf{การตรวจสอบความใช้ได้ของวิธีและการประเมินทางสถิติ:}
\begin{itemize}
    \item \textbf{การทดสอบ Paired Samples $t$-test:} ค่าความสูงต้นและมวลชีวภาพของกลุ่มที่ได้รับ PAW สูงกว่ากลุ่มควบคุมอย่างมีนัยสำคัญทางสถิติ ($t_{\text{stat}} = 18.42, p < 0.001$) สะท้อนว่าไนโตรเจนที่อยู่ในรูป $\ce{NO3-}$ จากพลาสมามีความสามารถในการดูดซึมเข้าสู่รากพืชได้อย่างมีประสิทธิภาพสูงกว่าปุ๋ยเคมีทั่วไป
    \item \textbf{การลดลงของเชื้อราก่อโรค:} ฤทธิ์ของกรดเพอร์ออกซีไนตรัส ($\ce{ONOOH}$) และ $\ce{H2O2}$ ในน้ำ PAW ส่งผลให้ปริมาณสปอร์ของเชื้อ *P. palmivora* ในดินรอบโคนต้นลดลงถึง 2.6-log (ร้อยละ 99.75) โดยไม่ส่งผลกระทบต่อจุลินทรีย์ตรึงไนโตรเจนและเชื้อไมคอร์ไรซาที่เป็นประโยชน์ในดิน
    \item \textbf{ความสอดคล้องแบลนด์-อัลต์แมน:} การเปรียบเทียบการวัดความเข้มข้นไนเตรตระหว่างเซนเซอร์เคมีไฟฟ้าแบบเรียลไทม์กับวิธีไอออนโครมาโตกราฟี (IC) ในห้องแล็บมาตรฐาน พบค่าความเอนเอียงเฉลี่ย $\bar{d} = -1.2 \text{ mg/L}$ ด้วยช่วงความสอดคล้อง $95\%$ Limits of Agreement อยู่ที่ $[-4.8, +2.4] \text{ mg/L}$ ผ่านเกณฑ์มาตรฐาน EURACHEM อย่างสมบูรณ์
\end{itemize}
\end{agribox}

\section{สรุปสาระสำคัญประจำบท}

น้ำกระตุ้นด้วยพลาสมา (PAW) เป็นนวัตกรรมปุ๋ยไนโตรเจนเหลวและสารอารักขาพืชทางชีวภาพที่ผลิตได้จากอากาศ น้ำ และพลังงานไฟฟ้า การถ่ายโอนมวลของอนุมูลออกไซด์ของไนโตรเจนตามกฎของเฮนรีเหนี่ยวนำให้เกิดสารประกอบไนโตรเจนที่มีประโยชน์ต่อการเจริญเติบโตของพืช ($\ce{NO3-}$ และ $\ce{NO2-}$) ขณะที่ปฏิกิริยาระหว่างกรดไนตรัสกับไฮโดรเจนเปอร์ออกไซด์ก่อให้เกิดกรดเพอร์ออกซีไนตรัส ($\ce{ONOOH}$) ซึ่งมีฤทธิ์ยับยั้งเชื้อราและแบคทีเรียก่อโรคพืชได้อย่างมีประสิทธิภาพโดยไร้สารเคมีตกค้าง การประยุกต์ใช้อัลกอริทึม NSGA-II ร่วมกับโมเดลโครงข่ายประสาทเทียมช่วยให้สามารถออกแบบระบบควบคุมสภาวะการผลิต PAW ให้ประหยัดพลังงานและตรงกับความต้องการของพืชแต่ละชนิดในแปลงปลูกได้อย่างแม่นยำ

\section{คำถามทบทวนท้ายบทที่ 5}

\begin{enumerate}
    \item จงอธิบายการถ่ายโอนมวลของอนุภาคจากพลาสมาลงสู่น้ำตามทฤษฎีสองฟิล์มและกฎของเฮนรี พร้อมระบุปัจจัยทางฟิสิกส์ที่มีผลต่อค่าสัมประสิทธิ์การถ่ายโอนมวล $k_{\text{L}}$
    \item จงเขียนสมการเคมีของการก่อตัวและการสลายตัวของกรดเพอร์ออกซีไนตรัส ($\ce{ONOOH}$) และอธิบายว่าเหตุใดสารชนิดนี้จึงมีฤทธิ์ทำลายเชื้อจุลินทรีย์สูงกว่าไฮโดรเจนเปอร์ออกไซด์เดี่ยวๆ
    \item เปรียบเทียบกระบวนการตรึงไนโตรเจนด้วยพลาสมากับกระบวนการฮาเบอร์-บอช ทั้งในมิติด้านอุณหพลศาสตร์ ผลกระทบต่อสิ่งแวดล้อม และความเป็นไปได้ในการกระจายการผลิตสู่ระดับฟาร์ม
    \item ระบบพลาสมาเครื่องหนึ่งใช้กำลังไฟฟ้า $150 \text{ W}$ ผลิตน้ำ PAW ปริมาตร $20 \text{ L}$ เป็นเวลา 45 นาที ได้ความเข้มข้นไนเตรตรวม $5.0 \text{ mM}$ จงคำนวณหาผลผลิตพลังงานในหน่วย g-N/kWh
    \item อธิบายหลักการของอัลกอริทึม NSGA-II ในการแก้ปัญหาความขัดแย้งระหว่างการเพิ่มอัตราการตรึงไนโตรเจนกับการประหยัดพลังงานไฟฟ้าในการผลิตปุ๋ย PAW
\end{enumerate}
